\documentclass[11pt]{article}
\usepackage[letterpaper,margin=0.78in]{geometry}
\usepackage{amsmath,amssymb,amsthm,mathtools}
\usepackage{booktabs,enumitem,microtype}
\usepackage[dvipsnames]{xcolor}
\usepackage[colorlinks=true,linkcolor=MidnightBlue,citecolor=MidnightBlue,urlcolor=MidnightBlue]{hyperref}
\setlist{nosep,leftmargin=*}
\newtheorem{theorem}{Theorem}
\newtheorem{proposition}[theorem]{Proposition}
\newtheorem{lemma}[theorem]{Lemma}
\newcommand{\cF}{\mathcal F}
\newcommand{\cC}{\mathcal C}
\newcommand{\R}{\mathbb R}
\title{Finite permutation cores for the nonunital product question\\
\large A rigorous partial result for arXiv:1705.08967}
\author{Autonomous solve-first investigation, lane 12}
\date{11 August 2026}

\begin{document}
\maketitle

\begin{abstract}
Niemiec and W\'ojcik ask whether a finite product of weakly (respectively
strongly) right amenable semigroups with topology (SGTs) remains right
amenable when the factors need not be unital.  We do not settle the general
question.  We give an exact finite-SGT reduction to a common permutation core,
use it to exclude every counterexample whose factors have order at most three,
and perform a complete self-product audit on 1,870 canonical order-four
extension types.  We also isolate why the usual iterated-mean proof breaks for
arbitrary SGTs.
\end{abstract}

\paragraph{The source question.}
An SGT is a semigroup equipped with an arbitrary topology; multiplication is
not assumed continuous.  For $f_s(x)=f(xs)$, the paper defines
$C_{\rm norm}(S)$ and $C_{\rm weak}(S)$ by requiring $f\in C_r(S)$ and
continuity of $s\mapsto f_s$ and $s\mapsto f_{st}$ in the norm or weak
topology, for every $t\in S$.  Immediately after Proposition 5.1(E), the
authors ask whether its unitality hypothesis can be dropped, even for a finite
family \cite[p.~14]{NW}.

\section{Exact finite reduction}

Let $S$ be finite.  Write $x\sim y$ when $x,y$ are in the same connected
component of the finite topology.  Components are clopen, so the continuous
real functions are exactly those constant on the $\sim$-classes.  Since every
function space below is finite dimensional, weak and norm continuity coincide.

Define $R$ to be the equivalence relation generated by the following pairs:
\begin{align*}
x&\mathrel R y &&(x\sim y),\\
xs&\mathrel R ys &&(x\sim y,\ s\in S),\\
xs&\mathrel R xs' &&(s\sim s',\ x\in S),\\
x(st)&\mathrel R x(s't) &&(s\sim s',\ x,t\in S).
\end{align*}

\begin{proposition}[finite function quotient]\label{prop:quotient}
$C_{\rm norm}(S)=C_{\rm weak}(S)$ is precisely the space of real functions
constant on the $R$-classes.  Moreover, $R$ is stable under every right
multiplication, so
\[
   \rho_s:S/R\longrightarrow S/R,\qquad [x]\longmapsto[xs]
\]
is well defined.
\end{proposition}

\begin{proof}
The first line is ordinary continuity.  The second is exactly the condition
$f_s\in C_b(S)$ for every $s$.  A map from a finite space to a Hausdorff vector
space is continuous exactly when it is constant on connected components; this
turns the two translation-orbit conditions into the third and fourth lines.
Thus the displayed relations are necessary and sufficient.  Stability under
right multiplication follows from associativity: the right translate of a
relation in lines one through four is respectively a relation in lines two,
two, four, and four.
\end{proof}

For maps $\rho_s$ on a finite set $Q$, call $K\subseteq Q$ a \emph{common
permutation core} if $K\ne\varnothing$ and every $\rho_s$ restricts to a
permutation of $K$.

\begin{theorem}[permutation-core criterion]\label{thm:core}
A finite SGT $S$ is weakly right amenable (equivalently, strongly right
amenable) if and only if the right action $\{\rho_s:s\in S\}$ on $S/R$ has a
common permutation core.
\end{theorem}

\begin{proof}
A mean on the finite-dimensional space of Proposition~\ref{prop:quotient} is a
probability vector $p$ on $Q=S/R$.  Right invariance says
$(\rho_s)_*p=p$ for every $s$.  For one deterministic map on a finite set, an
invariant probability gives zero mass to every transient point.  Hence each
$\rho_s$ permutes $\operatorname{supp}p$, proving necessity.  Conversely, the
uniform probability on a common permutation core is invariant under every
$\rho_s$ and therefore defines a right invariant mean.
\end{proof}

The largest possible core has a simple exact algorithm.  Start with $K_0=Q$.
Given $K_n$, retain only points that lie on a cycle of every $\rho_s$ contained
in $K_n$, and call the result $K_{n+1}$.  The process stabilizes in at most
$|Q|$ strict deletions.  Any common core survives every deletion, while a
nonempty fixed point is itself a common core.  This proves both correctness and
termination without numerical tolerance.

\section{Exhaustive and extended searches}

Every set partition occurs as the component partition of a finite topology
(take a disjoint union of indiscrete blocks), so enumerating partitions loses
no finite cases relevant to these function spaces.

\begin{center}
\begin{tabular}{@{}rrrr@{}}
\toprule
order & labelled semigroups & component partitions & amenable pairs\\
\midrule
1 & 1 & 1 & 1\\
2 & 8 & 2 & 15\\
3 & 113 & 5 & 534\\
\bottomrule
\end{tabular}
\end{center}

The enumeration tested every one of the 72,010 unordered pairs of nonunital
amenable candidates in these rows.  Every product had a nonempty permutation
core.  The core test was independently compared with a linear-program
feasibility test on all 582 factor cases and 500 sampled products; all 1,082
answers agreed.

To push past order three, a local associative extension enumerator began with
all 113 labelled order-three semigroups and exhausted the seven multiplication
entries involving a fourth element.  It produced 2,235 distinct labelled
order-four tables containing a three-element subsemigroup.  After testing all
15 component partitions and canonicalizing under all 24 relabellings, 1,870
amenable nonunital finite-SGT types remained.  \emph{Every one of their 1,870
self-products is amenable.}  In addition, 10,000 random cross-pairs from the
uncanonicalized catalog and 8,000 pairs from independently generated finite
transformation semigroups produced no counterexample.

These computations are exhaustive only where explicitly stated.  In
particular, the cross-pair order-four search and the larger transformation
search are evidence, not proofs.

\section{Why the general proof still fails}

For abstract amenable semigroups one writes
\[
 M(f)=m_S\bigl(s\mapsto m_T(t\mapsto f(s,t))\bigr).
\]
Without units this formula need not be defined on the paper's function spaces.
A slice $t\mapsto f(s,t)$ need not belong to $C_{\rm norm}(T)$ or
$C_{\rm weak}(T)$: continuity of $t\mapsto f(s,tv)$ cannot be obtained from a
product translate unless the first coordinate has an element that leaves $s$
fixed.  Extra right-factor anchors repair an individual slice, but after
averaging there is no reason for the outer translation-orbit map to remain norm
or weak continuous on a noncompact factor.

The finite criterion shows why the most obvious counterexamples also fail.
Product functions can distinguish extra ``indecomposable'' layers, but every
invariant probability assigns zero mass to layers transient under a right map.
To prove the finite product theorem one would still need to lift the product of
the two factor cores through the generally finer product function quotient.
The searches support such a lift, but a general proof was not obtained.

\paragraph{Scope.}
The unrestricted nonunital product question remains open in this packet.  The
rigorous contribution is Theorems~\ref{prop:quotient}--\ref{thm:core}, the
complete order-three exclusion, and a reproducible order-four self-product
audit.  Any future counterexample must evade all of these finite mechanisms;
the most plausible remaining setting is an infinite noncompact SGT where
averaging destroys translation-orbit continuity.

\begin{thebibliography}{9}
\bibitem{NW}
P. Niemiec and P. W\'ojcik,
\emph{Applications of amenable semigroups in operator theory},
Studia Math. 252 (2020), 27--48; arXiv:1705.08967.
\end{thebibliography}

\end{document}
